\documentclass[11pt, a4paper]{article}

% bfw-style.tex — gemeinsame Style-Definitionen für alle Handouts/Aufgabenhefte
% Voraussetzung: Kompilation mit xelatex (wegen fontspec)

\usepackage[ngerman]{babel}
\usepackage{geometry}
\geometry{a4paper, top=2.2cm, bottom=2.2cm, left=2.3cm, right=2.3cm}

\usepackage{fontspec}
% Fallback-Kette: Source Sans Pro -> Calibri -> Segoe UI -> Latin Modern Sans
% (Latin Modern ist immer mit MiKTeX vorhanden)
\IfFontExistsTF{Source Sans Pro}{%
  \setmainfont{Source Sans Pro}[Ligatures=TeX]%
}{\IfFontExistsTF{Calibri}{%
  \setmainfont{Calibri}[Ligatures=TeX]%
}{\IfFontExistsTF{Segoe UI}{%
  \setmainfont{Segoe UI}[Ligatures=TeX]%
}{%
  \setmainfont{Latin Modern Sans}[Ligatures=TeX]%
}}}
\IfFontExistsTF{Source Code Pro}{%
  \setmonofont{Source Code Pro}[Scale=0.92]%
}{\IfFontExistsTF{Consolas}{%
  \setmonofont{Consolas}[Scale=0.92]%
}{%
  \setmonofont{Latin Modern Mono}[Scale=0.92]%
}}

\usepackage{xcolor}
% Farbpalette: hell, freundlich, modern
\definecolor{bfwPrimary}{HTML}{0EA5A4}    % Petrol/Türkis - Primär
\definecolor{bfwPrimaryDark}{HTML}{0F766E} % dunkler Petrol für Headings
\definecolor{bfwAccent}{HTML}{F59E0B}     % Amber - Akzent
\definecolor{bfwSuccess}{HTML}{10B981}    % Grün - Erfolg / Lösung
\definecolor{bfwWarn}{HTML}{F97316}       % Orange - Achtung
\definecolor{bfwInk}{HTML}{0F172A}        % fast Schwarz
\definecolor{bfwInkSoft}{HTML}{475569}    % gedämpftes Grau
\definecolor{bfwBgSoft}{HTML}{F8FAFC}     % off-white
\definecolor{bfwBgTeal}{HTML}{ECFEFF}     % helles Türkis
\definecolor{bfwBgAmber}{HTML}{FEF3C7}    % helles Gelb
\definecolor{bfwBgRose}{HTML}{FFE4E6}     % helles Rosé für Eselsbrücken

\usepackage[colorlinks=true, linkcolor=bfwPrimaryDark, urlcolor=bfwPrimaryDark]{hyperref}
\usepackage{enumitem}
\setlist[itemize]{leftmargin=*, itemsep=2pt, topsep=2pt}
\setlist[enumerate]{leftmargin=*, itemsep=2pt, topsep=2pt}

\usepackage[most]{tcolorbox}
\usepackage{booktabs}
\usepackage{tikz}
\usetikzlibrary{decorations.pathmorphing, positioning, shapes.geometric, arrows.meta}

% Merkkästchen — wichtige Definition / Kernpunkt
\newtcolorbox{merksatz}[1][]{%
  enhanced, breakable,
  colback=bfwBgTeal,
  colframe=bfwPrimary,
  boxrule=0.6pt,
  arc=4pt,
  left=10pt, right=10pt, top=8pt, bottom=8pt,
  fonttitle=\bfseries\color{bfwPrimaryDark},
  title={\faLightbulb\ Merksatz},
  #1
}

% Eselsbrücke — Mnemoniks
\newtcolorbox{eselsbruecke}[1][]{%
  enhanced, breakable,
  colback=bfwBgRose,
  colframe=bfwAccent,
  boxrule=0.6pt,
  arc=4pt,
  left=10pt, right=10pt, top=8pt, bottom=8pt,
  fonttitle=\bfseries\color{bfwWarn},
  title={Eselsbrücke},
  #1
}

% Beispiel-Box
\newtcolorbox{beispiel}[1][]{%
  enhanced, breakable,
  colback=bfwBgSoft,
  colframe=bfwInkSoft,
  boxrule=0.4pt,
  arc=3pt,
  left=10pt, right=10pt, top=8pt, bottom=8pt,
  fonttitle=\bfseries\color{bfwInk},
  title={Beispiel},
  #1
}

% Lösungs-Box (für Aufgabenhefte)
\newtcolorbox{loesung}[1][]{%
  enhanced, breakable,
  colback=bfwBgAmber!40,
  colframe=bfwSuccess,
  boxrule=0.6pt,
  arc=3pt,
  left=10pt, right=10pt, top=8pt, bottom=8pt,
  fonttitle=\bfseries\color{bfwSuccess!70!black},
  title={Lösung},
  #1
}

% Glossar-Eintrag
\newcommand{\glossareintrag}[2]{%
  \par\noindent\textbf{\color{bfwPrimaryDark}#1} \enspace #2\par\smallskip
}

% Section-Stil — etwas farbiger als Standard
\usepackage{titlesec}
\titleformat{\section}
  {\Large\bfseries\color{bfwPrimaryDark}}
  {\thesection}{0.6em}{}
\titleformat{\subsection}
  {\large\bfseries\color{bfwPrimary}}
  {\thesubsection}{0.5em}{}
\titleformat{\subsubsection}
  {\normalsize\bfseries\color{bfwInk}}
  {\thesubsubsection}{0.4em}{}

% TikZ-Stil "skizzenhaft" — etwas weicher als Rough.js, aber stabil
\tikzset{
  roughbox/.style = {
    draw=bfwPrimaryDark, thick, fill=bfwBgTeal,
    rounded corners=4pt, inner sep=6pt
  },
  roughaccent/.style = {
    draw=bfwAccent, thick, fill=bfwBgAmber,
    rounded corners=4pt, inner sep=6pt
  },
  roughline/.style = {
    draw=bfwInkSoft, thick, -{Stealth[length=2.5mm]}
  }
}

% Bullet-Symbole (bewusst kein FontAwesome — vermeidet Font-Installations-Probleme)
\newcommand{\faLightbulb}{$\bullet$}
\newcommand{\faBookmark}{$\bullet$}

% Header/Footer
\usepackage{fancyhdr}
\pagestyle{fancy}
\fancyhf{}
\renewcommand{\headrulewidth}{0pt}
\fancyfoot[L]{\small\color{bfwInkSoft}\thepage}
\fancyfoot[R]{\small\color{bfwInkSoft} BFW M-064 Beschaffung im Büromanagement}


\title{\vspace*{-1.5cm}\Huge\bfseries\color{bfwPrimaryDark}Aufgabenheft Tag~6\\[0.4em]
  \large\color{bfwInkSoft}\textnormal{Terminplanung \& Terminmanagement --- Übungen im IHK-Klausurformat}}
\author{\textsf{Büroprozesse gestalten und Arbeitsvorgänge organisieren \textperiodcentered{} M-067}\\
\textsf{\small\copyright\ Can Siebert\,\textperiodcentered\,2026}}
\date{Selbstlernmaterial \textperiodcentered{} 06.07.2026}

\begin{document}
\maketitle
\thispagestyle{empty}

\begin{merksatz}[title={\bfseries So nutzen Sie dieses Aufgabenheft}]
Bearbeiten Sie alle Aufgaben zunächst \emph{ohne} in die Lösungen zu schauen. Schreiben
Sie Ihre Antworten in vollständigen Sätzen --- die IHK-Prüfung verlangt formulierte
Begründungen, nicht bloße Stichworte. Beachten Sie die \emph{Operatoren} (nennen,
erläutern, beurteilen, begründen) --- sie geben den geforderten Antwort-Tiefegrad vor.
Erst danach öffnen Sie den Lösungsteil ab Seite~\pageref{loesungen} und vergleichen.
\end{merksatz}

\tableofcontents
\vspace{1em}
\hrule

\section{Aufgaben}

\subsection*{Aufgabe~1 --- Was ist ein Termin? (10 Punkte)}

\noindent\textbf{\color{bfwAccent}YouTube-Vorbereitung:}\enspace \href{https://www.youtube.com/results?search_query=Terminplanung+Buero+Grundlagen}{Terminplanung --- Grundlagen} (\url{https://www.youtube.com/results?search_query=Terminplanung+Buero+Grundlagen})

\medskip
Sie sind im Sekretariat der \emph{Lindner \& Partner GmbH} tätig und sollen den
Terminkalender der Geschäftsführung pflegen.

\begin{enumerate}[label=(\alph*)]
  \item \textbf{Definieren} Sie den Begriff „Termin” und nennen Sie seine beiden Bestandteile. \hfill (3~P.)
  \item \textbf{Nennen} Sie drei typische betriebliche Termine. \hfill (3~P.)
  \item \textbf{Nennen} Sie die fünf W-Fragen, die vor dem Einplanen eines Termins zu klären sind. \hfill (2~P.)
  \item \textbf{Beurteilen} Sie die Aussage: \emph{„Bei einem Termin reicht das Datum, die Uhrzeit ist nebensächlich.”} \hfill (2~P.)
\end{enumerate}

\subsection*{Aufgabe~2 --- Terminarten (12 Punkte)}

\noindent\textbf{\color{bfwAccent}YouTube-Vorbereitung:}\enspace \href{https://www.youtube.com/results?search_query=feste+und+variable+Termine}{Feste und variable Termine} (\url{https://www.youtube.com/results?search_query=feste+und+variable+Termine})

\medskip
Termine lassen sich in feste und variable Termine unterscheiden.

\begin{enumerate}[label=(\alph*)]
  \item \textbf{Erläutern} Sie den Unterschied zwischen festen und variablen Terminen. \hfill (3~P.)
  \item \textbf{Ordnen} Sie folgende Termine \textbf{zu} (fest/variabel): Steuerzahlung, Besprechung, Betriebsferien, Geschäftsreise, Jubiläum, Geschäftsessen. \hfill (3~P.)
  \item \textbf{Erläutern} Sie den Begriff \emph{Serientermin} und nennen Sie zwei Beispiele. \hfill (3~P.)
  \item \textbf{Begründen} Sie, warum es sinnvoll ist, verschiedene Terminarten farblich zu kennzeichnen. \hfill (3~P.)
\end{enumerate}

\subsection*{Aufgabe~3 --- Terminüberwachung und Fristen (12 Punkte)}

\noindent\textbf{\color{bfwAccent}YouTube-Vorbereitung:}\enspace \href{https://www.youtube.com/results?search_query=Terminueberwachung+Wiedervorlage+Buero}{Terminüberwachung \& Wiedervorlage} (\url{https://www.youtube.com/results?search_query=Terminueberwachung+Wiedervorlage+Buero})

\medskip
Herr Fernandez aus der Buchhaltung findet eine zweite Mahnung des Lieferanten Opitz~AG
auf seinem Tisch.

\begin{enumerate}[label=(\alph*)]
  \item \textbf{Erläutern} Sie, was Terminüberwachung ist und wer dafür verantwortlich ist. \hfill (3~P.)
  \item \textbf{Nennen} Sie die Fälligkeit der Sozialabgaben sowie die Frist für die Lohnsteuer. \hfill (3~P.)
  \item \textbf{Nennen} Sie drei mögliche Folgen eines versäumten Zahlungstermins. \hfill (3~P.)
  \item \textbf{Erläutern} Sie, wie man sich verhalten sollte, wenn ein Termin nicht eingehalten werden kann. \hfill (3~P.)
\end{enumerate}

\subsection*{Aufgabe~4 --- Urlaubsplanung (15 Punkte)}

\noindent\textbf{\color{bfwAccent}YouTube-Vorbereitung:}\enspace \href{https://www.youtube.com/results?search_query=Urlaubsplanung+Paragraf+7+BUrlG}{Urlaubsplanung \& §\,7 BUrlG} (\url{https://www.youtube.com/results?search_query=Urlaubsplanung+Paragraf+7+BUrlG})

\medskip
In der Abteilung Beschaffung (Leitung: Frau Beckmann) vertreten sich Herr Anderson,
Frau Elling und Herr Lembke gegenseitig; inklusive Abteilungsleitung müssen immer zwei
Mitarbeiter anwesend sein. Im Mai darf wegen einer Produkteinführung kein Urlaub genommen
werden. Herr Lembke und Frau Elling möchten nur in den Schulferien Urlaub nehmen.

\begin{enumerate}[label=(\alph*)]
  \item \textbf{Erläutern} Sie, was Urlaubsplanung ist und warum sie eine besondere Form der Terminplanung darstellt. \hfill (4~P.)
  \item \textbf{Geben} Sie den Inhalt von §\,7 Abs.~1 und Abs.~2 BUrlG in eigenen Worten \textbf{wieder}. \hfill (5~P.)
  \item \textbf{Begründen} Sie, warum Frau Elling und Herr Lembke nicht gleichzeitig dieselbe Ferienwoche nehmen können, und \textbf{schlagen} Sie eine Lösung \textbf{vor}. \hfill (4~P.)
  \item \textbf{Beurteilen} Sie, welchen Nutzen eine frühzeitige Kalenderübersicht für die Urlaubsplanung hat. \hfill (2~P.)
\end{enumerate}

\subsection*{Aufgabe~5 --- Hilfsmittel und digitale Tools (12 Punkte)}

\noindent\textbf{\color{bfwAccent}YouTube-Vorbereitung:}\enspace \href{https://www.youtube.com/results?search_query=elektronischer+Terminkalender+Outlook}{Elektronischer Terminkalender} (\url{https://www.youtube.com/results?search_query=elektronischer+Terminkalender+Outlook})

\medskip
Ein Unternehmen will von einem Papierkalender auf eine digitale Lösung umsteigen.

\begin{enumerate}[label=(\alph*)]
  \item \textbf{Nennen} Sie drei klassische Hilfsmittel und drei Mindestfunktionen eines elektronischen Terminplanungsprogramms. \hfill (3~P.)
  \item \textbf{Nennen} Sie zwei Vorteile und einen Nachteil eines elektronischen Terminplaners gegenüber dem Papierkalender. \hfill (3~P.)
  \item \textbf{Erläutern} Sie, wofür sich ein Online-Tool wie \emph{Doodle} eignet und wann es gegenüber Outlook im Vorteil ist. \hfill (3~P.)
  \item \textbf{Beurteilen} Sie, welchen Nutzen geteilte Kalender und Terminbuchungs-Tools für die Zusammenarbeit im Team bieten. \hfill (3~P.)
\end{enumerate}

\subsection*{Aufgabe~6 --- Outlook und elektronische Besprechungen (10 Punkte)}

\noindent\textbf{\color{bfwAccent}YouTube-Vorbereitung:}\enspace \href{https://www.youtube.com/results?search_query=Outlook+Besprechung+planen+Teilnehmer+einladen}{Outlook --- Besprechung planen} (\url{https://www.youtube.com/results?search_query=Outlook+Besprechung+planen+Teilnehmer+einladen})

\medskip
\begin{enumerate}[label=(\alph*)]
  \item \textbf{Nennen} Sie die vier Terminarten von MS~Outlook und \textbf{ordnen} Sie je ein Beispiel \textbf{zu}. \hfill (4~P.)
  \item \textbf{Beschreiben} Sie in Stichpunkten, wie in Outlook eine Besprechungsanfrage erstellt wird. \hfill (3~P.)
  \item \textbf{Erläutern} Sie den Unterschied zwischen Videokonferenz und Telefonkonferenz und \textbf{nennen} Sie zwei Vorteile elektronischer Besprechungen. \hfill (3~P.)
\end{enumerate}

\vspace{1em}
\noindent\textbf{Punkte gesamt: 71}

\newpage

\section{Lösungen}\label{loesungen}

\subsection*{Lösung Aufgabe~1}
\begin{loesung}
\textbf{(a)} Ein Termin ist ein genau definierter Zeitpunkt, an dem etwas zu erledigen
ist oder ein Ereignis stattfindet. Er besteht aus den beiden Angaben \emph{Kalenderdatum}
und \emph{Uhrzeit}.

\textbf{(b)} Drei betriebliche Termine (Beispiele): Besprechungen, Geschäftsreisen,
Urlaub, Messen, Vertreterbesuche, Zahlungs- und Liefertermine.

\textbf{(c)} Die fünf W-Fragen: \emph{Warum} findet der Termin statt? \emph{Wer} nimmt
teil? \emph{Wann} (Datum, Uhrzeit)? \emph{Wo} (Ort, Raum)? \emph{Wie lange} dauert er?

\textbf{(d)} Die Aussage ist falsch. Ein Termin besteht aus Datum \emph{und} Uhrzeit;
ohne Uhrzeit ist ein Besprechungs- oder Lieferzeitpunkt nicht eindeutig planbar, und es
drohen Überschneidungen oder Versäumnisse.
\end{loesung}

\subsection*{Lösung Aufgabe~2}
\begin{loesung}
\textbf{(a)} Feste Termine wiederholen sich periodisch und werden möglichst frühzeitig
(teils zum Jahresanfang) eingetragen, um Terminüberschneidungen vorzubeugen. Variable
Termine werden sofort nach der Vereinbarung eingetragen; bei Bedarf erfolgt eine
Abstimmung mit bereits festliegenden Terminen.

\textbf{(b)} Zuordnung:
\begin{itemize}[itemsep=0pt]
  \item \emph{Fest}: Steuerzahlung, Betriebsferien, Jubiläum.
  \item \emph{Variabel}: Besprechung, Geschäftsreise, Geschäftsessen.
\end{itemize}

\textbf{(c)} Ein Serientermin ist ein Anlass, der sich nach einem Muster regelmäßig
wiederholt (täglich, wöchentlich, monatlich, jährlich). Beispiele: wöchentliches
Teammeeting, jährlicher Feiertag oder Geburtstag.

\textbf{(d)} Farbliche Kennzeichnung verschiedener Terminarten verschafft sofort optisch
einen Überblick: Auf einen Blick erkennt man, ob es sich z.\,B.\ um einen Kunden-,
Lieferanten- oder privaten Termin handelt --- das reduziert Fehler und Überschneidungen.
\end{loesung}

\subsection*{Lösung Aufgabe~3}
\begin{loesung}
\textbf{(a)} Terminüberwachung ist die Kontrolle der Einhaltung von Terminen. Sie wird
oft vom Sekretariat übernommen, gleichzeitig ist aber jeder selbst für die Einhaltung
seiner Termine verantwortlich. Hilfreich ist die Wiedervorlagemappe; am Ende des
Arbeitstages prüft man erledigte und für den Folgetag dringende Termine.

\textbf{(b)} Die Sozialabgaben (Arbeitnehmer- und Arbeitgeberanteil) sind am
\emph{drittletzten Bankarbeitstag des Monats} fällig; die einbehaltene Lohnsteuer ist
\emph{bis zum 10.\ des Folgemonats} an das Finanzamt abzuführen.

\textbf{(c)} Folgen eines versäumten Zahlungstermins: kein Skontoabzug mehr möglich,
Verzugszinsen, Mahngebühren, gestörtes Vertrauensverhältnis zum Lieferanten (drei
nennen).

\textbf{(d)} Ist eine Termineinhaltung nicht möglich, sollte der Termin sofort --- am
besten telefonisch und bei der zuständigen Person --- abgesagt und ein neuer Termin
vereinbart werden. Gut ist es, sofort einen Alternativtermin vorzuschlagen.
\end{loesung}

\subsection*{Lösung Aufgabe~4}
\begin{loesung}
\textbf{(a)} Bei der Urlaubsplanung werden die gewünschten oder genehmigten Urlaubstage
aller Mitarbeiter einer Abteilung in eine Kalenderübersicht eingetragen. Sie ist eine
besondere Form der Terminplanung, weil gleichzeitig die Rahmenbedingungen der Abteilung,
persönliche Wünsche und die gesetzliche Vorgabe des §\,7 BUrlG zu beachten sind.

\textbf{(b)} §\,7 Abs.~1: Bei der zeitlichen Festlegung sind die Urlaubswünsche des
Arbeitnehmers zu berücksichtigen --- es sei denn, dringende betriebliche Belange oder
unter sozialen Gesichtspunkten vorrangige Wünsche anderer Arbeitnehmer stehen entgegen.
Abs.~2: Der Urlaub ist grundsätzlich zusammenhängend zu gewähren; bei einem Anspruch von
mehr als zwölf Werktagen muss ein Urlaubsteil mindestens zwölf aufeinanderfolgende
Werktage umfassen.

\textbf{(c)} Problem: Beide werden für die gegenseitige Vertretung gebraucht, und es
müssen stets zwei Mitarbeiter anwesend sein --- sind beide gleichzeitig im Urlaub, ist
die Mindestbesetzung nicht gewährleistet. Lösung: Urlaub staffeln (z.\,B.\ erste vs.\
zweite Ferienhälfte), soziale Gesichtspunkte abwägen und frühzeitig in der
Kalenderübersicht abstimmen.

\textbf{(d)} Eine frühzeitige Kalenderübersicht macht Überschneidungen und
Engpässe sofort sichtbar, sichert die Mindestbesetzung und ermöglicht es, persönliche
Wünsche und gesetzliche Vorgaben rechtzeitig in Einklang zu bringen.
\end{loesung}

\subsection*{Lösung Aufgabe~5}
\begin{loesung}
\textbf{(a)} Klassische Hilfsmittel: Kalender, Planer, Terminmappen, Stecktafeln (drei
nennen). Mindestfunktionen eines elektronischen Programms: Kalenderfunktion,
Terminüberschneidungsprüfung, Adressbuch, Notizen, Erinnerungsfunktion, Passwortschutz
(drei nennen).

\textbf{(b)} Vorteile: Terminkalender können über ein Netzwerk verbunden und Termine
ausgetauscht, abgeglichen und gemeinsam bearbeitet werden; das Programm sucht auf Befehl
gemeinsame freie Termine (frei/belegt-Markierung); Erinnerung per optischer/akustischer
Meldung; mobiler Zugriff; leichtes Ändern/Verschieben/Löschen. Nachteil: Gefahr des
Datenverlustes bei Hard-/Softwareproblemen; Änderungen sind nicht mehr nachvollziehbar;
ohne Netzverbindung ist eine Synchronisation nötig.

\textbf{(c)} Doodle (oder Google Kalender) eignet sich zur Terminfindung mit externen
Partnern oder im privaten Bereich ohne MS~Outlook; mit einem selbsterklärenden
Assistenten lassen sich Terminabfragen erstellen. Vorteil gegenüber Outlook: Es ist
keine gemeinsame Outlook-/Exchange-Umgebung nötig --- ideal, wenn Teilnehmer aus
verschiedenen Organisationen abstimmen.

\textbf{(d)} Geteilte Kalender machen die Verfügbarkeit mehrerer Personen transparent und
erlauben die gemeinsame Pflege von Terminen. Terminbuchungs-Tools (z.\,B.\ Calendly,
Microsoft Bookings) lassen andere selbst freie Zeitfenster buchen --- das spart das
zeitaufwendige Hin und Her per E-Mail. Zu beachten ist der Datenschutz bei Cloud-Tools.
\end{loesung}

\subsection*{Lösung Aufgabe~6}
\begin{loesung}
\textbf{(a)} Die vier Outlook-Terminarten:
\begin{itemize}[itemsep=0pt]
  \item \emph{Ressourcen}: Räume, Fahrzeuge, Technik (z.\,B.\ Laptop, Beamer).
  \item \emph{Termine}: ohne weitere Personen und Ressource (z.\,B.\ Anruf, Arztbesuch).
  \item \emph{Besprechungen}: Einladung weiterer Personen plus Ressource (z.\,B.\ Besprechungsraum).
  \item \emph{Ereignisse}: 24 Stunden oder länger, ganztägig (z.\,B.\ Geburtstag, Feiertag, Urlaub).
\end{itemize}

\textbf{(b)} Besprechungsanfrage erstellen: Kalenderansicht öffnen $\rightarrow$ „Neue
Besprechung” $\rightarrow$ Titel und Ort eingeben $\rightarrow$ Datum/Beginn/Ende
(ggf.\ ganztägig) $\rightarrow$ erforderliche/optionale Teilnehmer aus dem Adressbuch
wählen $\rightarrow$ über „Terminplanung” freie Zeiten prüfen $\rightarrow$ „Senden”.

\textbf{(c)} Bei der Videokonferenz tauschen sich Gesprächspartner an verschiedenen
Orten per Video aus (spezielle Systeme + schneller Internetanschluss, auch gemeinsame
Dateiarbeit). Bei der Telefonkonferenz wird ein Telefonat mit mehr als drei Teilnehmern
geführt (genau drei = Dreierkonferenz). Vorteile elektronischer Besprechungen: keine
gleichzeitige Anwesenheit an einem Ort nötig; Einsparung von Flug-, Hotel- und
Verpflegungskosten sowie Zeit.
\end{loesung}

\end{document}
